By \(\text{thm:tcf-termination}\), the trace-independent base construction and the subsequent finite trace-annotation phase both terminate. By \(\text{lem:semialgebraic-fiber-dichotomy}\), the base formulas form a finite pairwise-disjoint partition of
\[
\mathcal P_{\mathrm{feas}}
=\{\mathbf p:\exists\theta\,\mathsf C_{\mathcal G}(\theta;\mathbf p)\}.
\tag{1}
\]
If \(P_{\mathrm{obs}}\) is compatible with the finite/full-positive class, the forward direction of \(\text{thm:response-fiber-equivalence}\) supplies \(\theta_{\mathcal M}\) satisfying \(\mathsf C_{\mathcal G}(\theta_{\mathcal M};P_{\mathrm{obs}})\) for any \(\mathcal M\in\mathfrak F(\mathcal G,P_{\mathrm{obs}})\). Hence its observed-cell vector belongs to (1) and lies in exactly one cell.

Fix that law and a coordinate \((a,y)\). Unfolding \(\text{def:terminal-qe-test}\) gives
\[
\begin{aligned}
\mathsf I_{a,y}(\mathcal G;P_{\mathrm{obs}})
&\Longleftrightarrow
\neg\exists\theta^0,\theta^1\bigl[
\mathsf C_{\mathcal G}(\theta^0;P_{\mathrm{obs}})
\wedge\mathsf C_{\mathcal G}(\theta^1;P_{\mathrm{obs}})\\
&\hspace{42mm}\wedge q_{a,y}(\theta^0)\neq q_{a,y}(\theta^1)
\bigr]\\
&\Longleftrightarrow q_{a,y}\text{ is constant on }
\{\theta:\mathsf C_{\mathcal G}(\theta;P_{\mathrm{obs}})\}.
\end{aligned}
\tag{2}
\]
The two directions of \(\text{thm:response-fiber-equivalence}\) are surjective in both directions between this polynomial fiber and \(\mathfrak F(\mathcal G,P_{\mathrm{obs}})\), and they preserve
\[
q_{a,y}(\theta)=Q_a^{\mathcal M}(y).
\tag{3}
\]
Therefore (2) holds if and only if \(\mathcal M\mapsto Q_a^{\mathcal M}(y)\) is constant on the finite/full-positive observed-law fiber.

Suppose first that (2) holds. The relation
\[
\mathsf V_{a,y}(P_{\mathrm{obs}},z)
\Longleftrightarrow
\exists\theta\,[\mathsf C_{\mathcal G}(\theta;P_{\mathrm{obs}})\wedge z=q_{a,y}(\theta)]
\tag{4}
\]
then has exactly one value of \(z\). If the cell carries a rational trace expression \(N/D\), the corrected \(\text{def:tcf-procedure}\) has verified on the entire cell that \(D\neq0\) and \(Dz=N\) whenever (4) holds. If no candidate passes that test, the stored functional is the unique-value graph (4) itself. In either case, let the stored value be \(f_{a,j}(P_{\mathrm{obs}})(y)\). For arbitrary \(\mathcal M\in\mathfrak F(\mathcal G,P_{\mathrm{obs}})\), the forward direction of response-fiber equivalence gives \(\theta_{\mathcal M}\) with
\[
\mathsf C_{\mathcal G}(\theta_{\mathcal M};P_{\mathrm{obs}}),
\qquad
q_{a,y}(\theta_{\mathcal M})=Q_a^{\mathcal M}(y).
\tag{5}
\]
Equations (4)--(5) and uniqueness yield the proved identifying equality
\[
f_{a,j}(P_{\mathrm{obs}})(y)=Q_a^{\mathcal M}(y)
\quad\text{for every }\mathcal M\in\mathfrak F(\mathcal G,P_{\mathrm{obs}}).
\tag{6}
\]

If (2) is false, the unique base cell carries \(\mathsf W_j\) by \(\text{lem:semialgebraic-fiber-dichotomy}\). The formula is satisfiable whether or not its real solutions have been given an effective encoding. Applying \(\text{prop:ambient-nonrecovery-transfer}\) gives strictly positive individual-edge-factorized models
\[
\mathcal M_0,\mathcal M_1\in\mathfrak F(\mathcal G,P_{\mathrm{obs}})
\subseteq\mathfrak M_{\mathrm{lm}}(\mathcal G),
\quad
P_{\mathcal M_0}(\mathcal O)=P_{\mathcal M_1}(\mathcal O)=P_{\mathrm{obs}},
\quad
Q_a^{\mathcal M_0}(y)\neq Q_a^{\mathcal M_1}(y).
\tag{7}
\]
Thus nonidentification in the ambient published class is an exact-free logical consequence. If \(\operatorname{Exact}(P_{\mathrm{obs}})\) is additionally supplied, \(\text{def:exact-observed-specialization}\) represents a satisfying pair over the real closure of the ordered field generated by \(\mathbf p\). Applying \(\text{thm:qe-failure-compiler}\) to that pair yields the models in (7) constructively and terminates with either a checked divergence record or an exhaustion tag. No colored-hedge conclusion is used or obtained.

Finally, suppose (6) holds for every \(y\in\mathcal X_Y\) at both \(a_0\) and \(a_1\). The outcome alphabet is finite by \(\text{ass:finite-alphabets}\). By \(\text{def:causal-query}\) and coordinatewise substitution from (6),
\[
\begin{aligned}
\tau_{s_Y}
&=\sum_{y\in\mathcal X_Y}s_Y(y)
\{Q_{a_1}^{\mathcal M}(y)-Q_{a_0}^{\mathcal M}(y)\}\\
&=\sum_{y\in\mathcal X_Y}s_Y(y)
\{f_{a_1,j}(P_{\mathrm{obs}})(y)-f_{a_0,j}(P_{\mathrm{obs}})(y)\}.
\end{aligned}
\tag{8}
\]
Equations (1)--(8) prove the pointwise semialgebraic characterization, the explicit subclass-relative identifying map, and the exact-free ambient obstruction on every negative cell.